\chapter{Previous Contact-Alignment Campaign}
\label{app:previous_campaign}

This appendix preserves the completed contact-alignment campaign that preceded
the axis-specific experiments. The measurements remain valid for the
configurations in which they were obtained and provide useful comparison data.
They are separated here because the surface plane and commanded tool
orientation were not independently parameterised: the controller used a
configured plane with \(a=0^\circ\), \(b=5^\circ\), whereas three manual
seatings estimated the physical plane as
\(a=-0.72\pm0.31^\circ\), \(b=14.99\pm1.91^\circ\). The resulting
\(10.01^\circ\) mismatch was approximately
\(-0.72^\circ\) about \(t_1\) and \(+9.99^\circ\) about \(t_2\).

\section{Test Matrix}

The campaign contains \(92\) recorded runs. The \(88\) contact runs form the
five cases in \cref{tab:previous_campaign_cases}. One additional mode-0 hold
record is usable only as a null-space diagnostic; three further hold records
are excluded because their general log is missing or their Cartesian task
drift exceeds the registered limit.

\begin{table}[H]
\centering
\small
\caption{Previous contact-alignment cases.}
\label{tab:previous_campaign_cases}
\begin{tabular}{@{}
>{\centering\arraybackslash}p{0.07\textwidth}
>{\raggedright\arraybackslash}p{0.22\textwidth}
>{\raggedright\arraybackslash}p{0.36\textwidth}
>{\raggedright\arraybackslash}p{0.25\textwidth}@{}}
\toprule
Case & Test & Design & Purpose \\
\midrule
A &
Set-up settling &
\(4\), \(8\), and \(12\,\mathrm{s}\); one run per duration &
Select a settled evaluation time. \\
B &
Formulation equivalence &
Decoupled command (\(n=5\)) and zero-shift coupled command (\(n=3\)) &
Check agreement of equivalent impedance forms. \\
C &
Nominal repeatability &
Five nominal runs with a physical reset between runs &
Establish the run-to-run interpretation scale. \\
D &
Paired tipping stiffness &
\(K_{R,t_1}=K_{R,t_2}=5,\ 15,\ 50\,\mathrm{Nm/rad}\);
five runs per setting &
Measure the response to a common tangent-axis stiffness. \\
E &
Compliance-centre position &
Edge check (\(n=3\)); configured-normal sweep (\(n=24\));
\(t_1\) sweep (\(n=18\)); \(t_2\) sweep (\(n=12\)) &
Measure the response to the three pole coordinates. \\
\bottomrule
\end{tabular}
\end{table}

\section{Validation Results}

The \(4\,\mathrm{s}\) set-up ended when the commanded press ramp completed and
still contained the load overshoot. The estimated normal load relaxed by
approximately \(4.6\,\mathrm{s}\). Re-evaluation of the longer traces at
\(5\,\mathrm{s}\) produced the residual changes in
\cref{tab:previous_five_second_check}; all are below the campaign's
repeatability-based interpretation threshold.

\begin{figure}[H]
    \centering
    \includegraphics[width=\textwidth]{G2_equilibrium.pdf}
    \caption{Previous set-up settling comparison.}
    \label{fig:previous_settling}
\end{figure}

\begin{table}[H]
\centering
\caption{Previous residual changes after \(5\,\mathrm{s}\).}
\label{tab:previous_five_second_check}
\begin{tabular}{@{}lccc@{}}
\toprule
Trace & Tip difference & Alignment difference & Load difference \\
\midrule
8-second run  & \(0.006^\circ\) & \(0.017^\circ\) & \(0.09\,\mathrm{N}\) \\
12-second run & \(0.003^\circ\) & \(0.002^\circ\) & \(0.12\,\mathrm{N}\) \\
\bottomrule
\end{tabular}
\end{table}

The decoupled and zero-shift coupled commands differed by \(0.079^\circ\) in
final tip angle, or \(0.34\) times the larger group standard deviation. No
discrepancy exceeding the measured repeatability was detected. Five nominal
repetitions gave the dispersions in
\cref{tab:previous_repeatability}.

\begin{table}[H]
\centering
\caption{Previous nominal repeatability (\(n=5\)).}
\label{tab:previous_repeatability}
\begin{tabular}{@{}lccc@{}}
\toprule
Quantity & Mean & \(\sigma\) & \(2\sigma\) \\
\midrule
Final tip angle (\(^\circ\)) & 1.474 & 0.194 & 0.388 \\
Configured-normal change (\(^\circ\)) & \(-0.714\) & 0.294 & 0.588 \\
Estimated normal load (N) & 59.752 & 0.484 & 0.969 \\
Edge travel (mm) & 21.146 & 0.589 & 1.178 \\
\bottomrule
\end{tabular}
\end{table}

\section{Paired Stiffness Results}

Both tangent-axis gains were changed together in Case D. The data therefore
characterise the common setting \(K_{R,t_1}=K_{R,t_2}\); they do not identify
the separate sensitivities to the two axes. The imposed mismatch was
predominantly about \(t_2\), which further limits an axis-specific
interpretation.

\begin{table}[H]
\centering
\caption{Previous paired-stiffness results (\(n=5\)).}
\label{tab:previous_stiffness_results}
\begin{tabular}{@{}lccc@{}}
\toprule
\(K_{R,t_1}=K_{R,t_2}\) (Nm/rad)
& Tip (\(^\circ\))
& Estimated normal load (N)
& Improvement (\(^\circ\)) \\
\midrule
5  & \(2.539\pm0.132\) & \(60.27\pm0.27\) & \(+2.38\pm0.06\) \\
15 & \(1.841\pm0.066\) & \(60.77\pm0.22\) & \(+1.59\pm0.04\) \\
50 & \(1.280\pm0.032\) & \(60.53\pm0.23\) & \(+0.79\pm0.04\) \\
\bottomrule
\end{tabular}
\end{table}

\begin{figure}[H]
    \centering
    \includegraphics[width=\textwidth]{A2_stiffness_sweep.pdf}
    \caption{Previous paired tangent-axis stiffness sweep.}
    \label{fig:previous_stiffness_sweep}
\end{figure}

Reducing both stiffnesses from \(50\) to \(5\,\mathrm{Nm/rad}\) increased the
mean physical-plane alignment improvement from \(0.79^\circ\) to
\(2.38^\circ\). The mean model-estimated normal load remained between
\(60.27\) and \(60.77\,\mathrm{N}\).

\section{Compliance-Centre Results}

The position campaign contains \(57\) runs. An additive quadratic model fitted
to the \(54\) coordinate-sweep runs accounts for \(R^2=0.962\) of the response
using \(t_1\), \(t_1^2\), \(t_2\), and \(t_2^2\). It predicts a local maximum
near \(t_1=+62\,\mathrm{mm}\), \(t_2=+30\,\mathrm{mm}\), but this coordinate
pair was not tested directly. Three contact-edge runs were held out as a
consistency check.

\begin{figure}[H]
    \centering
    \includegraphics[width=\textwidth]{B_pole_component.pdf}
    \caption{Previous response by compliance-centre coordinate.}
    \label{fig:previous_pole_components}
\end{figure}

\begin{figure}[H]
    \centering
    \includegraphics[width=\textwidth]{B_pole_surface.pdf}
    \caption{Previous fitted in-plane response.}
    \label{fig:previous_pole_surface}
\end{figure}

The response was non-monotonic. Unsuitable in-plane placements produced
alignment improvements as low as \(-10.31^\circ\). Adding the physical normal
coordinate to the in-plane model increased \(R^2\) only from \(0.962\) to
\(0.966\) over the investigated range. This is specific to the tested contact
geometry and does not establish that the normal coordinate is generally
irrelevant.

\section{Interpretation Limits}

The physical plane was estimated by seating the tool face rather than from
three calibrated probe points. Absolute residual angles therefore inherit an
angular spread of approximately \(\pm1.9^\circ\), although the within-run
improvement changes by less than \(0.1^\circ\) across the three estimates.
The workpiece fixture moved under load, the tool loosened during two repeated
runs, and the reported contact load comes from the robot's model-estimated
external wrench rather than an independent force/torque sensor.

The compliance-centre samples form a cross rather than a complete
two-dimensional grid, so the fitted model contains no \(t_1t_2\) interaction.
The predicted local maximum and its transfer to another plane orientation,
tool, surface, or robot configuration remain unverified.
